% chapters/requirements/rf_export.tex
% RF-14: Exportación y búsqueda
% ============================================================================

\item \textbf{Exportación y búsqueda.} \label{rf:export}
El sistema debe proporcionar capacidades de búsqueda avanzada para los gestores,
permitiendo la exploración jerárquica del histórico del sistema (cursos $\rightarrow$
asignaturas $\rightarrow$ exámenes $\rightarrow$ instancias) y la exportación de datos en
formatos estándar. La exportación de notas en formato compatible con sistemas de actas
institucionales permite la integración con los flujos administrativos existentes.

\begin{functional}

    % RF-14.1 ---------------------------------------------------------------
    \item \textbf{Exportación de notas en formato institucional.}
    \label{rf:export:institutional}
    El sistema debe permitir al coordinador exportar las calificaciones de un examen en un
    formato \acs{csv} compatible con los sistemas institucionales de actas.

    \textbf{Entrada:} petición GET al recurso de exportación de notas del examen con
    parámetro de formato.

    \textbf{Proceso:} el sistema genera un fichero \acs{csv} con la estructura adecuada
    (columnas para identificador del alumno, nombre, grupo, calificación). Se incluyen las
    calificaciones de todas las instancias del examen cuyo estado sea \texttt{PUBLISHED} o
    posterior. El formato de columnas es configurable por organización para adaptarse a las
    especificaciones de cada institución.

    \textbf{Salida:} respuesta con el fichero \acs{csv} como descarga.

    % RF-14.2 ---------------------------------------------------------------
    \item \textbf{Datos completos de calificaciones.} \label{rf:export:raw}
    El sistema debe permitir obtener los datos completos de calificaciones de un examen en
    formato estructurado.

    \textbf{Entrada:} petición GET al recurso de calificaciones del examen.

    \textbf{Proceso:} el sistema devuelve los datos completos de cada instancia: alumno
    (\acs{nia}, nombre), grupo, modelo, calificación de cada problema y nota total. Estos
    datos permiten al \dfn{frontend} generar estadísticas (por grupo, modelo, problema,
    etc.) sin necesidad de lógica adicional en el \dfn{backend}.

    \textbf{Salida:} respuesta con los datos en formato \acs{json}.

    % RF-14.3 ---------------------------------------------------------------
    \item \textbf{Búsqueda jerárquica por gestor.} \label{rf:export:hierarchical}
    El sistema debe permitir al gestor explorar el histórico de su organización de forma
    jerárquica: cursos $\rightarrow$ asignaturas $\rightarrow$ exámenes $\rightarrow$
    instancias.

    \textbf{Entrada:} petición GET a los recursos de asignaturas, exámenes o instancias con
    parámetros de filtrado y el identificador del curso (activo o archivado).

    \textbf{Proceso:} el sistema permite la navegación jerárquica dentro de la organización
    del gestor: al consultar una asignatura se listan sus exámenes; al consultar un examen
    se listan sus instancias. Se soporta la navegación entre cursos archivados. Los
    resultados son de solo lectura para datos de cursos archivados.

    \textbf{Salida:} respuesta con los datos del nivel consultado, con enlaces a los niveles
    inferiores.

    % RF-14.4 ---------------------------------------------------------------
    \item \textbf{Filtros avanzados sobre entidades.} \label{rf:export:filters}
    El sistema debe soportar filtros avanzados en los listados de todas las entidades
    principales.

    \textbf{Entrada:} parámetros de filtrado en las peticiones GET a los recursos de
    listado. Los filtros disponibles varían según la entidad: usuarios (nombre, estado),
    asignaturas (nombre, código, cuatrimestre), exámenes (nombre, estado), instancias
    (alumno, grupo, modelo, estado, \texttt{has\_issues}).

    \textbf{Proceso:} el sistema aplica los filtros de forma combinable (AND) y devuelve los
    resultados paginados. Se soporta búsqueda parcial por texto en campos de nombre.

    \textbf{Salida:} respuesta con los resultados paginados y los metadatos de paginación.

    % RF-14.5 ---------------------------------------------------------------
    \item \textbf{Acceso a datos de cursos archivados.} \label{rf:export:archived}
    El sistema debe garantizar que los datos de cursos archivados permanecen accesibles en
    modo de solo lectura para los gestores de la organización.

    \textbf{Entrada:} cualquier petición a recursos de un curso archivado.

    \textbf{Proceso:} el sistema identifica que el curso del recurso solicitado está
    archivado. Las operaciones de lectura se permiten con normalidad. Cualquier operación de
    escritura se rechaza. Pueden coexistir múltiples cursos archivados accesibles
    simultáneamente.

    \textbf{Salida:} datos en modo lectura, o error \acs{http}~403 si se intenta una
    escritura.

\end{functional}
